\section{Recover a sextet from an input character}
Decoding reverses the alphabet mapping. An uppercase character \(x\) represents
\(x-65\); lowercase represents \(x-71\); a digit represents \(x+4\).
\texttt{+} and \texttt{/} represent 62 and 63. Every other input octet is invalid
except a padding character handled at its permitted final position.

The scalar classifier first tries the uppercase interval. Subtract its lower bound,
then compare that unsigned difference with the interval width. A character below
\texttt{A} produces a wrapped large unsigned difference and fails the comparison.
A character above \texttt{Z} produces a difference at least 26 and also fails.

\Needspace{10\baselineskip}
\MARSCodeLabel{Recognize an ASCII interval with one unsigned range test.}
\begin{minted}[samepage=true]{mars-idea}
b64_digit:
    daddiu $t0, $a0, -65
    sltiu $t1, $t0, 26
    bnez $t1, digit_direct
    nop
\end{minted}

Use the same reasoning for lowercase and digits, then test the two final symbols.
The classifier returns a sextet or minus one; the validator checks that return before
using it. A zero sextet is a valid result for \texttt{A}, so a zero/nonzero test cannot
distinguish success from failure.

The alternative implementation uses a 256-byte decode table, with minus one in invalid
positions. The input byte is zero-extended before indexing; the table result is loaded
signed so its minus-one marker remains distinguishable. Both strategies must agree for
all 256 possible octets, including zero, spaces, newlines and bytes at least 128.
